\documentclass[11pt,a4paper]{article}
\usepackage[utf8]{inputenc}
\usepackage[T1]{fontenc}
\usepackage[french]{babel}
\usepackage{amsmath,amssymb}
\usepackage{geometry}
\usepackage{booktabs,longtable,array}
\geometry{margin=2.2cm}
\setlength{\parskip}{0.4em}
\newcommand{\R}{\mathbb{R}}

\title{Théorie de la loi de commande hybride\\ \large Linéarisation par retour d'état + commande prédictive (MPC) contrainte\\ appliquée à la GADA (DFIG) chaotique et au MPPT}
\author{}
\date{}

\begin{document}
\maketitle

\section{Modèle à commander}

\subsection{Modèle physique de la GADA}
Avec l'orientation du flux statorique sur l'axe $q$ ($\psi_{qs}=\psi_s$, $\psi_{ds}=0$), $u_{ds}=-V_s$, $u_{qs}=0$ et $\psi_s=V_s/\omega_s$, le vecteur d'état est
\[
x=\begin{bmatrix}x_1\\x_2\\x_3\end{bmatrix}=\begin{bmatrix}i_{dr}\\ i_{qr}\\ \omega_r\end{bmatrix},
\]
et la dynamique s'écrit
\begin{equation}
\left\{\begin{aligned}
\dot x_1 &= a_1x_1+(\omega_s-x_3)x_2+a_2x_3+c_1+u_1\\
\dot x_2 &= -(\omega_s-x_3)x_1+a_1x_2+c_2+u_2\\
\dot x_3 &= a_6x_1-a_7x_3-a_8
\end{aligned}\right.
\label{eq:modele}
\end{equation}
avec $\sigma=1-L_m^2/(L_sL_r)$ et
\begin{align*}
a_1&=-\frac{R_sL_m^2}{\sigma L_s^2L_r}-\frac{R_r}{\sigma L_r}, &
a_2&=-\frac{L_m\psi_s}{\sigma L_rL_s}, &
a_3&=-\frac{L_m}{\sigma L_sL_r}, &
a_4&=\frac{1}{\sigma L_r},\\
a_5&=\frac{R_sL_m\psi_s}{\sigma L_s^2L_r}, &
a_6&=\frac{3p^2L_m\psi_s}{2JL_s}, &
a_7&=\frac{D}{J}, &
a_8&=\frac{p}{J}T_L,\\
c_1&=a_3u_{ds}, & c_2&=a_5+a_3u_{qs}. & &
\end{align*}
Le couple de charge est l'opposé du couple aérodynamique ramené à l'arbre de la génératrice, $T_L=-T_a$, avec
\[
T_a=\frac{\tfrac12\rho\pi R^2C_p(\lambda)v^3}{\Omega_m},\qquad \Omega_m=\frac{\omega_r}{p},\qquad
\lambda=\frac{R\,\Omega_m}{G\,v},
\]
\[
C_p=0.5176\Big(\frac{116}{\lambda_i}-5\Big)e^{-21/\lambda_i}+0.0068\lambda,\qquad \frac1{\lambda_i}=\frac1\lambda-0.035 .
\]

Les seules entrées de commande sont $u_1=a_4\,\Delta u_{dr}$ et $u_2=a_4\,\Delta u_{qr}$, c'est-à-dire les tensions rotoriques imposées par le convertisseur côté rotor. L'équation de vitesse n'a \textbf{pas d'actionneur} : la vitesse est commandée à travers le couple électromagnétique $T_e\propto i_{dr}$, donc à travers $i_{dr}$ (structure en cascade).

On écrit les deux premières équations sous forme compacte :
\begin{equation}
\dot x_{1,2} = F_{1,2}(x,t)+u_{1,2} .
\label{eq:compact}
\end{equation}

\subsection{Modèle adimensionnel (étude du chaos)}
Avec $\tau=-1/a_1$, $\tilde t=t/\tau$, $k=a_7/(\tau a_6)$, $z_1=x_1/k$, $z_2=x_2/k$ et $z_3=\tau(x_3-\omega_s)$ :
\begin{equation}
\left\{\begin{aligned}
z_1'&=-z_1-z_3z_2+\gamma z_3+\varepsilon_1\\
z_2'&=-z_2+z_3z_1+\varepsilon_2\\
z_3'&=s(z_1-z_3)+\varepsilon_3
\end{aligned}\right.
\qquad
\gamma=\frac{a_2a_6\tau}{a_7}<0,\quad s=\frac{\tau D}{J}.
\label{eq:adim}
\end{equation}
Ce modèle a la même structure, avec $u_1,u_2$ sur les deux premières équations et aucune entrée sur la troisième. La loi de commande ci-dessous s'applique donc aux deux cas.

\section{Étape 0 : boucle externe de vitesse (cascade)}

Comme $T_e=\tfrac32p\frac{L_m}{L_s}\psi_s\,i_{dr}$, on commande la vitesse en imposant la référence de $i_{dr}$. Soit $e_\omega=\omega_r-\omega_{r,ref}$ et $I_\omega=\int e_\omega\,dt$. On choisit
\begin{equation}
\boxed{i^\ast_{dr}=\frac{\dot\omega_{r,ref}+a_7\omega_r+\hat a_8(v,\omega_r)-\lambda e_\omega-\lambda_I I_\omega}{a_6}}
\label{eq:vitesse}
\end{equation}
En remplaçant dans la troisième équation du modèle :
\[
\dot e_\omega=-\lambda e_\omega-\lambda_I I_\omega+a_6\,(i_{dr}-i^\ast_{dr}).
\]
\begin{itemize}
\item Si $i_{dr}=i^\ast_{dr}$, l'erreur de vitesse obéit à $\ddot I_\omega+\lambda\dot I_\omega+\lambda_I I_\omega=0$ : système linéaire stable dont les pôles sont les racines de $p^2+\lambda p+\lambda_I$. Avec $\lambda=10$ et $\lambda_I=25$, on obtient un pôle double en $-5$ (amortissement critique, sans dépassement).
\item Sinon, c'est un système stable excité par l'erreur de la boucle interne $i_{dr}-i^\ast_{dr}$. Il est donc stable entrée-état (ISS), et $e_\omega\to0$ dès que la boucle interne converge (argument de cascade).
\item La dérivée $\dot i^\ast_{dr}$, nécessaire à la linéarisation de la boucle interne, est calculée analytiquement à partir du modèle :
\[
\dot i^\ast_{dr}=\frac{\ddot\omega_{r,ref}+a_7\dot\omega_r+\dot{\hat a}_8-\lambda(\dot\omega_r-\dot\omega_{r,ref})-\lambda_I e_\omega}{a_6},
\]
où $\dot\omega_r=F_3(x)$ et $\dot{\hat a}_8$ s'obtient par dérivation en chaîne de $T_a(v,\omega_r)$.
\end{itemize}
La référence de la boucle interne est $x_d=[i^\ast_{dr},\ i_{qr,ref}]^T$.

\section{Étape 1 : linéarisation par retour d'état (boucle interne)}

Soit $e=x_{1,2}-x_d$ l'erreur de poursuite des courants. On choisit
\begin{equation}
\boxed{u_{1,2} = -\hat F_{1,2}(x,t)+\dot x_d+v}
\label{eq:FL}
\end{equation}
où $\hat F_{1,2}$ est le modèle nominal (paramètres nominaux) et $v\in\R^2$ une nouvelle entrée. En remplaçant \eqref{eq:FL} dans \eqref{eq:compact} :
\begin{equation}
\dot e = \dot x_{1,2}-\dot x_d = F_{1,2}-\hat F_{1,2}+v = v+\Delta(x,t),\qquad \Delta=F_{1,2}-\hat F_{1,2} .
\label{eq:erreur}
\end{equation}
\begin{itemize}
\item \textbf{Cas nominal ($\Delta=0$).} On obtient deux intégrateurs découplés $\dot e_i=v_i$. Les termes bilinéaires $(\omega_s-x_3)x_{1,2}$, responsables du chaos, sont compensés exactement.
\item \textbf{Avec erreur de modèle ($\Delta\neq0$).} $\Delta$ agit comme une perturbation. Une loi purement proportionnelle $v=-k_e e$ laisse alors une erreur statique, ce qui justifie l'action intégrale de la section suivante.
\end{itemize}

\section{Étape 2 : modèle d'erreur augmenté et discrétisé}

Pour chaque canal $i\in\{1,2\}$, on ajoute l'intégrale de l'erreur $\sigma_i=\int e_i\,dt$ :
\[
\xi_i=\begin{bmatrix}e_i\\ \sigma_i\end{bmatrix},\qquad
\dot\xi_i=\underbrace{\begin{bmatrix}0&0\\1&0\end{bmatrix}}_{A_c}\xi_i+\underbrace{\begin{bmatrix}1\\0\end{bmatrix}}_{B_c}v_i .
\]
On discrétise exactement, avec un bloqueur d'ordre zéro de période $T_s$ :
\[
A=e^{A_cT_s}=I+A_cT_s=\begin{bmatrix}1&0\\T_s&1\end{bmatrix},\qquad
B=\int_0^{T_s}e^{A_c\eta}B_c\,d\eta=\begin{bmatrix}T_s\\T_s^2/2\end{bmatrix}.
\]
Le développement est exact car $A_c^2=0$. On obtient :
\begin{equation}
\boxed{\xi_i(k+1)=A\,\xi_i(k)+B\,v_i(k)}
\label{eq:discret}
\end{equation}

\section{Étape 3 : commande prédictive (MPC) contrainte}

\subsection{Prédiction}
Sur un horizon de $N$ pas, avec $V=[v(k|k),\dots,v(k+N-1|k)]^T$, la prédiction s'écrit
\[
\xi(k+j|k)=A^j\xi(k)+\sum_{l=0}^{j-1}A^{j-1-l}B\,v(k+l|k),\qquad j=1,\dots,N .
\]
Sous forme matricielle :
\[
\Xi=\Phi\,\xi(k)+\Gamma V,\quad
\Phi=\begin{bmatrix}A\\A^2\\ \vdots\\ A^N\end{bmatrix},\quad
\Gamma=\begin{bmatrix}B&0&\cdots&0\\ AB&B&\cdots&0\\ \vdots&&\ddots&\\ A^{N-1}B&A^{N-2}B&\cdots&B\end{bmatrix}.
\]

\subsection{Critère}
\begin{equation}
J(V)=\sum_{j=1}^{N}\|\xi(k+j|k)\|_Q^2+\|\xi(k+N|k)\|_P^2+r\sum_{j=0}^{N-1}v(k+j|k)^2 ,
\label{eq:cout}
\end{equation}
où :
\begin{itemize}
\item $Q=\mathrm{diag}(q_e,q_I)>0$ pondère l'erreur et son intégrale ;
\item $r>0$ pénalise l'effort de commande ;
\item $P$ est le poids terminal donné par la LMI (section~\ref{sec:lmi}).
\end{itemize}
Avec $Q=W^TW$ (Cholesky) et $\bar W=\mathrm{blkdiag}(W,\dots,W,\mathrm{chol}(Q+P))$, le critère s'écrit comme un problème de moindres carrés :
\[
J=\big\|\,\underbrace{\begin{bmatrix}\bar W\Gamma\\ \sqrt r I_N\end{bmatrix}}_{W_G}V+\underbrace{\begin{bmatrix}\bar W\Phi\\0\end{bmatrix}}_{W_P}\xi(k)\big\|^2
= V^TH V+2g^TV+\text{cte},\quad H=W_G^TW_G,\ g=W_G^TW_P\,\xi(k).
\]

\subsection{Contraintes}
L'actionneur est borné : $|u_i|\le\bar u_i$. D'après \eqref{eq:FL}, $u_i=h_i+v_i$ avec $h=-\hat F(x,t)+\dot x_d$, d'où
\begin{equation}
\underline v_i=-\bar u_i-h_i\ \le\ v_i\ \le\ \bar u_i-h_i=\overline v_i .
\label{eq:contraintes}
\end{equation}
La valeur de $h_i$ est gelée sur l'horizon (elle est mesurée à l'instant $k$).

\subsection{Problème QP et horizon glissant}
\begin{equation}
\boxed{V^\ast=\arg\min_V\ V^THV+2g^TV\quad\text{s.c.}\quad \underline v\le V\le\overline v},\qquad v(k)=V^\ast_1 .
\label{eq:qp}
\end{equation}
Seul le premier élément $V^\ast_1$ est appliqué, puis le problème est résolu à nouveau à l'instant $k+1$ (horizon glissant). Comme les contraintes sont des bornes simples (boîte), le QP est résolu par une méthode d'ensemble actif. Le résultat est identique à celui de \texttt{quadprog}.

\subsection{Anti-emballement (anti-windup)}
L'intégrale n'est mise à jour que si la commande n'est pas saturée :
\[
\sigma_i(k+1)=\sigma_i(k)+T_s\,e_i(k)\quad\text{si } |u_i(k)|<\bar u_i,\qquad \sigma_i(k+1)=\sigma_i(k)\ \text{sinon}.
\]

\section{Étape 4 : poids terminal et gain par LMI (stabilité)}\label{sec:lmi}

\textbf{Théorème.} S'il existe $X=X^T>0$ et $Y$ tels que
\begin{equation}
\begin{bmatrix}
X & (AX+BY)^T & XQ^{1/2} & \sqrt r\,Y^T\\
AX+BY & X & 0 & 0\\
Q^{1/2}X & 0 & I & 0\\
\sqrt r\,Y & 0 & 0 & 1
\end{bmatrix}\succeq0,
\label{eq:lmi}
\end{equation}
alors, avec $P=X^{-1}$ et $K=YX^{-1}$,
\begin{equation}
(A+BK)^TP(A+BK)-P+Q+rK^TK\preceq0 .
\label{eq:lyap}
\end{equation}

\textbf{Démonstration.}
\begin{enumerate}
\item Le complément de Schur de \eqref{eq:lmi} par rapport au bloc diagonal $\mathrm{diag}(X,I,1)$ donne
\[
X-(AX+BY)^TX^{-1}(AX+BY)-XQX-rY^TY\succeq0 .
\]
\item On multiplie à gauche et à droite par $P=X^{-1}$ et on pose $Y=KX$ ; on obtient \eqref{eq:lyap}.
\item $V_f(\xi)=\xi^TP\xi$ est donc une fonction de Lyapunov pour la loi locale $v=K\xi$ :
\[
V_f(\xi^+)-V_f(\xi)\le-\|\xi\|_Q^2-r\,v^2 .
\]
\item L'ensemble $\Omega=\{\xi^TP\xi\le\alpha\}$, avec $\alpha$ tel que $K\xi$ respecte les bornes dans $\Omega$, est invariant.
\item Par l'argument classique coût terminal / ensemble terminal (Mayne \emph{et al.}, 2000), le coût optimal $V_N^\ast$ décroît :
\[
V_N^\ast(k+1)-V_N^\ast(k)\le-\|\xi(k)\|_Q^2-r\,v(k)^2 .
\]
Le MPC est donc récursivement faisable, et $\xi=0$ (c'est-à-dire $x\to x_d$) est asymptotiquement stable dans le cas nominal. \hfill$\square$
\end{enumerate}

\textbf{Calcul pratique.} La solution de volume maximal de \eqref{eq:lmi} rend \eqref{eq:lyap} égale à zéro : $P$ est la solution de l'équation de Riccati discrète (DARE)
\[
P=A^TPA-A^TPB(r+B^TPB)^{-1}B^TPA+Q,\qquad K=-(r+B^TPB)^{-1}B^TPA .
\]
Dans le programme, $P$ est calculée par l'algorithme de doublement, et la LMI \eqref{eq:lmi} est vérifiée numériquement. La plus petite valeur propre obtenue vaut environ $10^{-16}$, soit $\approx0$ : la LMI est satisfaite à la limite, comme prévu.

\textbf{Conséquences.}
\begin{itemize}
\item Sans contrainte active, le MPC donne exactement $v=K\xi=k_e e+k_I\int e$, c'est-à-dire un PI appliqué au système linéarisé.
\item Un seul gain $K=[k_e,\ k_I]$ est commun aux deux canaux.
\end{itemize}

\section{Loi de commande complète et lois de comparaison}
\begin{center}
\begin{tabular}{lp{6.2cm}p{5.2cm}}
\toprule
Loi & Expression & Rôle\\
\midrule
FL--MPC (proposée) & $i^\ast_{dr}$ par \eqref{eq:vitesse} ; $u_{1,2}=\mathrm{sat}\big(-\hat F_{1,2}+\dot x_d+v_{MPC}\big)$ & compensation + optimisation sous contraintes + intégrale\\
FL seule & même boucle de vitesse ; $u_{1,2}=\mathrm{sat}\big(-\hat F_{1,2}+\dot x_d-k_e e\big)$ & loi de la version soumise (sans intégrale)\\
PI en cascade & $i^\ast_{dr}=-(\lambda e_\omega+\lambda_I I_\omega)/a_6$ ; $u_{1,2}=\mathrm{sat}\big(-k_e e-k_I\!\int e\big)$ & commande vectorielle classique, sans modèle, avec anti-windup\\
\bottomrule
\end{tabular}
\end{center}
Les trois lois agissent uniquement sur $u_1,u_2$ et utilisent les mêmes gains $k_e$, $k_I$ (LMI) et $\lambda$, $\lambda_I$. La comparaison est donc équitable : elle isole l'apport de la linéarisation, puis celui du MPC.

\section{Références de poursuite (MPPT)}
\begin{align*}
\omega_{r,ref}&=pG\lambda_{opt}\,\bar v/R && \text{(vitesse optimale, vent filtré } \bar v)\\
i_{dr,ref}&=\big(\dot\omega_{r,ref}+a_7\omega_{r,ref}+(p/J)\,T_L(v,\omega_{r,ref})\big)/a_6 && \text{(équilibre des couples)}\\
i_{qr,ref}&=\psi_s/L_m && (Q_s=0,\ \text{facteur de puissance unitaire})\\
P_s&=-\tfrac32V_s\frac{L_m}{L_s}\,i_{dr} && \text{(puissance active statorique)}
\end{align*}
Le vent est modélisé par
\[
v(t)=V_{mean}+\sum_{k=1}^3A_k\sin(2\pi f_kt+\phi_k)+v_{turb}(t),
\]
où la turbulence est un bruit blanc filtré au premier ordre :
\[
v_{turb}(k)=v_{turb}(k-1)+\frac{\Delta t}{\tau_t}\Big(-v_{turb}(k-1)+\sigma_t\sqrt{2\tau_t/\Delta t}\;w(k)\Big),\qquad w(k)\sim\mathcal N(0,1).
\]

\section{Cible de la suppression du chaos}
La cible $z^\ast$ est choisie comme un point d'équilibre réalisable. On impose $\omega_r=1.2\,\omega_s$ (fonctionnement hypersynchrone, vitesse positive), soit
\[
z_3^\ast=0.2\,\omega_s\tau=2.02 .
\]
Deux conditions complètent le choix :
\begin{itemize}
\item $z_3'=0$ (pas d'actionneur sur la vitesse, $\varepsilon_3=0$) donne $z_1^\ast=z_3^\ast$ ;
\item $z_1'=0$ avec $u_1=0$ donne $z_2^\ast=(-z_3^\ast+\gamma z_3^\ast+\varepsilon_1)/z_3^\ast=-2.00$.
\end{itemize}
Cet équilibre est maintenu par $u_2$ seule :
\[
u_2^\ast=-(-z_2^\ast+z_3^\ast z_1^\ast+\varepsilon_2)=78.9 ,
\]
ce qui compense l'excitation anormale $\varepsilon_2$.

\section{Justification des valeurs du programme}

\small
\begin{longtable}{>{\raggedright}p{3.3cm}p{2.3cm}p{9.6cm}}
\toprule
Paramètre & Valeur & Justification\\
\midrule
\endhead
\multicolumn{3}{l}{\textbf{Génératrice (tableau I)}}\\
$R_s$, $R_r$ & 2.6, 2.9 m$\Omega$ & Ordres de grandeur d'une GADA de classe 2 MW, grandeurs ramenées au stator. À vérifier ou remplacer par une fiche technique réelle si disponible.\\
$L_m$ ; $L_s=L_r$ & 2.5 ; 2.587 mH & Inductances de fuite de 87 $\mu$H, d'où $L_s,L_r>L_m$ et $\sigma=0.0661\in(0,1)$ : condition physique exigée par les relecteurs.\\
$p$ & 2 & Machine à 4 pôles, cas usuel des GADA de forte puissance.\\
$V_{ll}$, $f_s$ & 690 V, 50 Hz & Tension standard des éoliennes MW ; réseau algérien et européen à 50 Hz. On en déduit $\psi_s=1.793$ Wb et $V_s=563$ V (crête par phase).\\
$J$ & 127 kg$\cdot$m$^2$ & Inertie totale (rotor et turbine) ramenée côté génératrice, ordre de grandeur réaliste pour 2 MW.\\
$D$ & $10^{-3}$ N$\cdot$m$\cdot$s/rad & Frottement visqueux faible, cas normal.\\
$\tau=-1/a_1$ & 32.1 ms & Calculée à partir des paramètres ; c'est la constante de temps électrique rotorique.\\
\midrule
\multicolumn{3}{l}{\textbf{Turbine}}\\
$\rho$ & 1.225 kg/m$^3$ & Air standard (15 °C, niveau de la mer).\\
$R$ & 40 m & Diamètre de 80 m, typique d'une éolienne de 2 MW ; $P_a\approx1.5$ MW à 10 m/s.\\
$G$ & 80 & Choisi pour que la vitesse synchrone soit atteinte vers 9.7 m/s : la vitesse reste entre $0.82\omega_s$ et $1.12\omega_s$, dans la plage usuelle de glissement ($\pm30$ \%).\\
$\lambda_{opt}$, $C_{p,max}$ & 8.1 ; 0.48 & Maximum de la courbe $C_p(\lambda)$ utilisée (formule classique de Heier).\\
\midrule
\multicolumn{3}{l}{\textbf{Chaos (modèle adimensionnel)}}\\
$\gamma$, $s$ & $-1$ ; 8 & $\gamma<0$ est imposé par la physique ($\sigma>0$). $s=8$ place le système dans la zone chaotique (mode mécanique plus rapide que $\tau$).\\
$\varepsilon=(0,-85,0)$ & & D'après le diagramme de bifurcation, le chaos apparaît pour $\varepsilon_2<-15$. La valeur $-85$ est au cœur de la zone chaotique, loin des fenêtres périodiques ($-50.5$, $-41.8$, $-28.5$).\\
$z_0$ & (0.3, 0.6, 1) & Condition initiale quelconque, conservée de la version soumise.\\
$\Delta t$ (RK4) & $10^{-3}$ & Pas petit devant le temps caractéristique du chaos ($1/\lambda_{max}\approx0.7$).\\
$T_{trans}$, $T_{lyap}$ & 200 ; 500 & Élimination du transitoire, puis moyenne longue pour faire converger les exposants de Lyapunov. Contrôle : la somme des exposants vaut $-(2+s)=-10$.\\
Grille $\varepsilon_2$ & 401 points & Pas de 0.25, suffisant pour localiser le seuil ($\approx-15$) et les fenêtres périodiques.\\
\midrule
\multicolumn{3}{l}{\textbf{Suppression du chaos (partie 2)}}\\
$T_s$ & 0.01 & En unités de $\tau$, soit 0.32 ms réelles, période typique d'une commande de convertisseur (quelques kHz).\\
$N$ & 15 & Horizon de $0.15\tau$, plus long que la constante de temps rapide en boucle fermée ($\approx0.01$).\\
$Q=\mathrm{diag}(1,10)$, $r=10^{-4}$ & & On obtient $k_e=-63.9$ et $k_I=-192.4$. Les pôles en boucle fermée valent 0.38 et 0.969, ce qui donne une convergence en $\approx0.3$ unité de temps. Le poids 10 sur l'intégrale élimine rapidement l'erreur statique.\\
$\bar u$ & 120 & Supérieur à $u_2^\ast=78.9$ (l'équilibre reste atteignable, marge de 50 \%) mais inférieur aux commandes demandées au départ du chaos : les contraintes sont actives, ce qui teste le MPC.\\
$\tilde t_{on}$ & 6 & Laisse voir clairement le régime chaotique avant l'activation.\\
Désadaptation & 0, $\pm20$ \% & Incertitude paramétrique usuelle (résistances avec la température, inductances avec la saturation, inertie mal connue).\\
\midrule
\multicolumn{3}{l}{\textbf{Vent}}\\
$V_{mean}$ & 9 m/s & Sous la vitesse nominale ($\approx11$ m/s), donc zone MPPT (sans action du calage).\\
$A_k$, $f_k$ & (1.3, 0.7, 0.4), (0.05, 0.13, 0.31) Hz & Variations lentes, de période 3 à 20 s.\\
$\sigma_t$, $\tau_t$ & 0.9 m/s ; 0.5 s & Intensité de turbulence $\approx10$ \% ; turbulence rapide, de l'ordre de la seconde.\\
Filtre MPPT & 2 s & La turbine (grande inertie) ne peut pas suivre les rafales rapides : la vitesse de référence suit le vent filtré.\\
\midrule
\multicolumn{3}{l}{\textbf{Poursuite MPPT (partie 3)}}\\
$T_s$ & 1 ms & Commande à 1 kHz, soit $T_s\ll\tau=32$ ms.\\
$N$ & 10 & Horizon de 10 ms $\approx\tau/3$.\\
$Q=\mathrm{diag}(1,3)$, $r=10^{-6}$ & & On obtient $k_e=-619$ et $k_I=-1070$. La bande passante de la boucle de courant est $\approx620$ rad/s ($\approx100$ Hz), valeur classique pour une boucle de courant de GADA.\\
$|\Delta u_{dr,qr}|$ & 200 V & $\approx0.35\,V_s$ : tension rotorique d'un convertisseur dimensionné pour un glissement de $\pm30$ \% ($0.3\times563=169$ V), plus une marge.\\
$\lambda$, $\lambda_I$ & 10 s$^{-1}$, 25 s$^{-2}$ & Boucle de vitesse : pôle double en $-5$ rad/s (constante de temps 0.2 s), environ 100 fois plus lente que la boucle de courant ($\approx620$ rad/s). Cette séparation des échelles de temps justifie la structure en cascade. Mêmes valeurs (normalisées) pour le cas chaotique.\\
Mise à l'échelle & (1 kA, 1 kA) & Améliore le conditionnement numérique du QP. Elle ne change pas le gain $K$ (invariance d'échelle).\\
M1 / M2 & $\pm20$ \% ($R$, $J$), $\mp10$ \% ($L_m$) & Même logique d'incertitude. Les inductances de fuite sont conservées, donc $L_s,L_r>L_m$ reste vrai.\\
Pas d'intégration & $T_s/4$, $T_s/5$ & Précision de RK4 face à la rotation à $\omega_s=314$ rad/s (période de 20 ms).\\
Métriques & $t\ge5$ s & Exclut le transitoire initial (mise en vitesse, $\approx0.9$ s).\\
\bottomrule
\end{longtable}
\normalsize

\paragraph{Remarque honnête sur le choix des réglages.}
\begin{itemize}
\item Les valeurs de $Q$, $r$ et $N$ sont des choix de conception : ce sont des compromis entre rapidité et effort de commande, et non des valeurs uniques.
\item Pour que la comparaison soit juste, les trois lois (FL--MPC, FL seule, PI) utilisent exactement le même gain $K$ issu de la LMI.
\item Les paramètres de la machine sont représentatifs d'une classe 2 MW, et non tirés de la fiche technique d'une machine précise.
\end{itemize}

\end{document}
